\section{Hearing Harmonics}\label{sec:harmonics}

A spectrum connects a sound's timbre to its partials. Try this with an
oscillator that offers sine, triangle, saw, and square outputs at one pitch.

\paragraph{Starting point}
Load \textbf{Harmonics}. Set all input gains to $0\,\mathrm{dB}$, Length
4096, Hop $30\,\mathrm{ms}$, and Y Scale to Log 60dB. Keep Average and
Smooth off, Slope at zero, and the view at $0$--$5\,\mathrm{kHz}$.
Tune the oscillator near $220\,\mathrm{Hz}$; exact tuning is not required.

\begin{enumerate}
  \item Patch sine to red and saw to green. Look for a shared lowest peak,
  the fundamental. The saw should have peaks near integer multiples of it:
  $440$, $660$, $880\,\mathrm{Hz}$, and so on.
  \item Patch triangle to blue and a square with 50\% pulse width to yellow.
  Compare peaks near $3f_0$ and $5f_0$. An ideal triangle's upper harmonics
  fall faster than a square's; both lack even harmonics.
  \item Change the pulse width. Even harmonics can now appear, and some
  harmonics weaken or disappear. Listen while watching the pattern change.
  \item Increase Smooth, then return it to None. The broad tonal shape is
  easier to compare when smoothed, but individual partials merge.
\end{enumerate}

\begin{figure}[!htp]
\centering
% Ideal relative partial amplitudes, independent of module UI and FFT results.
\begin{tikzpicture}[x=1cm,y=1cm,font=\small]
\foreach \origin/\title/\mode in {0/Saw/1,4.5/Square/2,9/Triangle/3} {
  \begin{scope}[shift={(\origin,0)}]
    \node[anchor=west,font=\small\bfseries,text=fourierAccent] at (0,2.25) {\title};
    \draw[->,panelMuted] (0,0)--(3.85,0);
    \draw[->,panelMuted] (0,0)--(0,1.85);
    \foreach \n in {1,...,7} {
      \pgfmathsetmacro{\amp}{\mode==1 ? 1/\n : (mod(\n,2)==0 ? 0 : (\mode==2 ? 1/\n : 1/(\n*\n)))}
      \draw[fourierAccent,line width=1.2pt] (\n*.45,0)--(\n*.45,\amp*1.65);
      \fill[fourierAccent] (\n*.45,\amp*1.65) circle[radius=.035];
    }
    \node[below] at (.45,0) {$f_0$};
    \node[below] at (1.35,0) {$3f_0$};
    \node[below] at (2.25,0) {$5f_0$};
    \node[below] at (3.15,0) {$7f_0$};
  \end{scope}
}
\end{tikzpicture}

\caption{Ideal harmonic amplitudes relative to each waveform's fundamental.
Stems indicate discrete partials, not the broadened peaks of a finite FFT.
Real oscillators may intentionally depart from these shapes.}
\end{figure}

\paragraph{What to learn}
Peak spacing reveals harmonic structure; relative heights help explain
brightness. Equal peak voltage at the oscillator outputs does not guarantee
equal fundamental height or perceived loudness. To compare only spectral
shape, temporarily adjust analyzer input gains to align the fundamentals.
Return to equal gains before making a level comparison. A missing peak may
also lie below the selected Y Scale or be hidden beneath another trace.
